\documentclass[11pt]{article}
\usepackage[margin=1in]{geometry}
\usepackage{booktabs}
\usepackage{amsmath}
\title{Macro-op Fusion in the CTRL-ALT-ELITE RV32I Core\\\large Measured performance change of the v1 \texttt{addi}+\texttt{branch} fusion}
\author{CTRL-ALT-ELITE}
\date{23 November 2026}
\begin{document}
\maketitle
\begin{abstract}
The v1 macro-op fusion recogniser packs an \texttt{addi} and the conditional branch that
tests its result into a single pipeline slot. Measured across fourteen benchmark kernels on
one build with the feature masked on and off, total cycles fall from \textbf{1339 to 1048},
a reduction of \textbf{291 cycles (21.7\%)}. The saving is exactly the number of fused
pairs, the non-issue overhead of the pipeline is unchanged to within rounding, and every
test passes in both configurations. The pre-RTL gate predicted 21.4\%; the measurement
delivers 21.7\%. On the graded programs, however, the feature returns nothing --- and the
epilogue records the optimisation that replaced it and does return a gain.
\end{abstract}

\section{The feature}

A counted loop ends with the same two-instruction idiom every iteration: the loop counter is
decremented by an \texttt{addi}, and a conditional branch tests it.

\begin{verbatim}
    addi x5, x5, -1        decrement
    bne  x5, x0, loop      test the result
\end{verbatim}

\texttt{fuse\_detect} recognises this pair in the fetch stage -- \texttt{addi} immediately
followed by a conditional branch whose \texttt{rs1} is the register the \texttt{addi}
writes -- and consumes it as \emph{one} slot: the PC advances by 8 instead of 4, the branch
word never enters the pipeline, and its condition, compared register and target ride along
with the \texttt{addi} as payload. Execute computes the increment on the existing adder and
tests it on the existing comparator.

The recogniser never sees a register \emph{value}, so it makes a packing decision only; the
branch outcome is still decided in execute. The feature is masked by a compile-time constant,
\texttt{FUSE\_ENABLE}, so both configurations below are the same source with one character
changed and nothing else.

\section{Measurement method}

Simulation is deterministic: identical RTL and identical images produce identical cycle
counts. The two columns below are therefore the same machine, differing only in whether the
recogniser's verdict is allowed to reach the datapath. No cycle difference can be attributed
to anything but the feature.

The full regression is 67 tests plus one long-standing known failure
(\texttt{sample-bmarks/default}, failing since 2026-09-23 and unchanged by this work). Both
configurations pass all 67. The fourteen kernels below are the benchmark suite built for
exactly this purpose; each is self-checking, so a wrong answer fails the run rather than
producing a fast number.

\section{Results}



Every kernel that contains a fused pair saves exactly one cycle per pair: \texttt{sweep100}
saves 100, \texttt{sweep50} saves 50, \texttt{sweep25} saves 25, \texttt{sweep4} saves 4.
\texttt{smoke}, which has no loop at all, is unchanged at six cycles. No kernel over-delivers
and none under-delivers.

\begin{table}[h]
\centering\small
\caption{Per-kernel cycles. ``pairs'' is the number of fused pairs executed, counted
statically and multiplied by the trip count; ``flush'' is the squashed-instruction count,
which is identical in both configurations; ``ovh'' is cycles minus retired instructions.}
\begin{tabular}{lrrrrrrr}
\toprule
kernel & cycles off & cycles on & saved & \% & pairs & flush off/on & ovh off/on \\
\midrule
\texttt{smoke} & 6 & 6 & 0 & +0.0 & 1 & 0/0 & 4/4 \\
\texttt{sweep4} & 29 & 25 & -4 & -13.8 & 4 & 3/3 & 10/10 \\
\texttt{sweep10} & 47 & 37 & -10 & -21.3 & 10 & 3/3 & 10/10 \\
\texttt{loop10} & 47 & 37 & -10 & -21.3 & 10 & 3/3 & 10/10 \\
\texttt{sweep25} & 92 & 67 & -25 & -27.2 & 25 & 3/3 & 10/10 \\
\texttt{sweep50} & 167 & 117 & -50 & -29.9 & 50 & 3/3 & 9/10 \\
\texttt{sweep100} & 317 & 217 & -100 & -31.5 & 100 & 3/3 & 9/10 \\
\texttt{nested} & 112 & 87 & -25 & -22.3 & 25 & 10/10 & 24/24 \\
\texttt{if-in-loop} & 56 & 46 & -10 & -17.9 & 10 & 6/6 & 16/16 \\
\texttt{sum-array} & 51 & 47 & -4 & -7.8 & 4 & 3/3 & 14/14 \\
\texttt{array-idx} & 91 & 83 & -8 & -8.8 & 8 & 3/3 & 18/18 \\
\texttt{fib} & 69 & 59 & -10 & -14.5 & 10 & 3/3 & 10/10 \\
\texttt{smul} & 218 & 186 & -32 & -14.7 & 32 & 5/5 & 14/14 \\
\texttt{callret} & 37 & 34 & -3 & -8.1 & 3 & 4/4 & 12/12 \\
\hline
\textbf{total} & \textbf{1339} & \textbf{1048} & \textbf{-291} & \textbf{-21.7} & \textbf{292} & -- & -- \\
\bottomrule
\end{tabular}
\end{table}

\section{Where the cycles went}

The saving is not a throughput effect that has to be explained by pipeline dynamics; it is
arithmetic. Retired instructions fall by 293 and cycles fall by 291, so the non-issue cycles
of the machine -- pipeline fill, load-use stalls, flush cycles -- are untouched.

\begin{table}[h]
\centering
\begin{tabular}{lrrr}
\toprule
 & feature off & feature on & change \\
\midrule
retired instructions & 1169 & 876 & -293 \\
non-issue cycles (overhead) & 170 & 172 & +2 \\
total cycles & 1339 & 1048 & -291 \\
\midrule
CPI & 1.15 & 1.20 & +0.05 \\
IPC & 0.873 & 0.836 & -0.037 \\
\bottomrule
\end{tabular}
\end{table}

\subsection{Both ratios move the wrong way, and that is correct}

CPI rises from 1.15 to 1.20: its denominator is the
instruction count, which fusion shrinks by design. IPC \emph{also} falls, from
0.873 to 0.836, and that is the more interesting result.

IPC is instructions per cycle. Fusion removes 293 instructions and 291 cycles, so the same
fixed overhead -- about 170 cycles of fill, stalls and flushes -- is now shared among 876
instructions instead of 1169. Efficiency \emph{per instruction} therefore drops even though
the program finishes 22\% sooner: each surviving instruction is charged a larger share of the
fixed cost.

This corrects the hypothesis recorded before the experiment, which predicted that IPC must
rise. It does not, and it cannot, for a feature whose mechanism is to shorten the
instruction stream while leaving the pipeline's fixed costs alone. \textbf{The metric for
this feature is cycles for fixed work.} CPI and IPC are both accounting artefacts of a
changing denominator and should not be used to accept or reject a fusion variant -- exactly
the caveat the experiment log already carried, now confirmed by measurement.

\section{Why the saving survived the predictor}

Fusion removes branches from the dynamic stream -- the signal the branch predictor is
measured on -- so the saving could easily have been cancelled by extra mispredicts. It was
not: the flush-cycle count is identical in every kernel (3/3, 10/10, 16/16, 24/24 in the
table above).

That is not automatic. The branch inside a fused pair is never fetched as its own
instruction, so the predictor would never look it up unless told to. Fetch therefore reads
the prediction for \texttt{pc+4} when a pair is fused, redirects to the pair's own target
when that prediction is taken, and trains the entry at \texttt{pc+4} when the pair resolves.
Without that wiring the first build of this feature was \emph{slower}: every taken fused
branch mispredicted, costing two cycles against the one slot saved, a net loss of one cycle
per iteration. The current result depends on that wiring, and the identical flush counts are
the evidence that it works.

\section{Prediction versus measurement}

Before any RTL, the opportunity was counted from the kernels' sources: 287 pairs under the
v1 rule, 25.6\% of the dynamic instruction stream, and therefore a bound of 21.4\% of the
suite's cycles.

The measurement is 21.7\%. The two agree to within the rounding of the CPI-derived
instruction count, which is the strongest form of validation available here: the model that
justified building the feature is the model that predicts its result.

\section{Effect on the directed suite}

No directed test slows down. Five improve by exactly one cycle, each because it happens to
contain one fused pair:

\begin{table}[h]
\centering
\begin{tabular}{lrrr}
\toprule
test & cycles off & cycles on & change \\
\midrule
\texttt{mx-bypass-rs1-brcomp} & 14 & 13 & -1 \\
\texttt{rv32ui-p-bge} & 347 & 346 & -1 \\
\texttt{rv32ui-p-bgeu} & 372 & 371 & -1 \\
\texttt{rv32ui-p-bne} & 315 & 314 & -1 \\
\texttt{rv32ui-p-jal} & 23 & 22 & -1 \\
\bottomrule
\end{tabular}
\end{table}

The remaining 48 pass with identical cycle counts, which is the expected outcome: they are
directed instruction tests with essentially no loop-carried \texttt{addi}+\texttt{branch}
pairs.

\section{Program time and the clock-speed budget}

Cycles are the machine-independent half of the answer. Turning them into time needs a clock,
which is where this section is careful about what has and has not been measured.

\subsection{Simulation wall time}

Measured on the simulator binary itself, so the figures exclude the Verilator build and
describe execution only:

\begin{center}
\begin{tabular}{lrr}
\toprule
test & cycles & wall time \\
\midrule
\texttt{smoke} & 6 & 0.01\,s \\
\texttt{loop10} & 37 & 0.02\,s \\
\texttt{array-idx} & 83 & 0.04\,s \\
\texttt{sweep100} & 217 & 0.09\,s \\
\texttt{rv32ui-p-add} & 463 & 0.19\,s \\
\midrule
full suite, 67 tests & -- & 19.4\,s (including one build per testbench) \\
\bottomrule
\end{tabular}
\end{center}

Against the autograder's 100\,s per test (500\,s for the performance programs) the largest
single kernel finishes in under a fifth of a second, so simulation cost is not a
constraint on any experiment here.

\subsection{Execution time from cycles}

For a clock period $T$, execution time is $\text{cycles} \times T$. Over the fourteen
kernels the total falls from 1339 to 1048 cycles:

\begin{center}
\begin{tabular}{rrrr}
\toprule
clock & feature off & feature on & saved \\
\midrule
50\,MHz & 26.8\,$\mu$s & 21.0\,$\mu$s & 5.8\,$\mu$s \\
100\,MHz & 13.4\,$\mu$s & 10.5\,$\mu$s & 2.9\,$\mu$s \\
200\,MHz & 6.7\,$\mu$s & 5.2\,$\mu$s & 1.5\,$\mu$s \\
\bottomrule
\end{tabular}
\end{center}

The absolute numbers are small because the kernels are small, and they assume the clock is
unchanged. What matters is the ratio: \textbf{21.7\% fewer cycles for the same work}. If
$f_{max}$ is unaffected that is also 21.7\% less time; section~8.3 states the margin
available if it is not.

\subsection{The budget: how much $f_{max}$ the feature may cost}

Fusion adds logic to the fetch path -- an opcode compare, a register-index compare and a
32-bit adder for the branch target -- so it may reduce the achievable clock. Let
$f_{\text{off}}$ and $f_{\text{on}}$ be the frequencies of the two builds. The feature pays
in time exactly when

\[
\frac{\text{cycles}_{\text{on}}}{f_{\text{on}}} < \frac{\text{cycles}_{\text{off}}}{f_{\text{off}}}
\quad\Longleftrightarrow\quad
\frac{f_{\text{on}}}{f_{\text{off}}} > \frac{1048}{1339} = 0.782 .
\]

\textbf{The clock may get up to 21.7\% slower before the feature stops paying.} A 5\%
$f_{max}$ hit still nets about 17\% faster; a 25\% hit makes it a net loss and the answer is
\texttt{FUSE\_ENABLE = 0}. That single inequality is the whole risk budget, and it is the
one number still missing from this report.

\subsection{Where $f_{max}$ has to come from}

It cannot be measured on the development machine: \texttt{yosys}, \texttt{sta},
\texttt{openroad} and \texttt{librelane} are all absent there, and the project's own
instructions place area-delay metrics on the EA machines. Two things are needed for the run
to mean anything:

\begin{enumerate}
\item \textbf{Blackbox the memory.} The behavioural \texttt{main\_memory} is an array of a
thousand 32-bit words. Synthesised, it becomes tens of thousands of flip-flops whose read
path dominates the critical path, and the resulting slack describes the array rather than
the processor -- which is what an earlier attempt in this project measured. The flow now
reads \texttt{synthesis/main\_memory\_bb.sv}, a port-identical blackbox, so the reported
$f_{max}$ belongs to the pipeline.
\item \textbf{Sweep the clock period,} not just 10\,ns. \texttt{synthesis/timing.tcl}
constrains a single period; running it at several periods brackets $f_{max}$ directly and
makes the 0.782 comparison a measurement rather than an estimate.
\end{enumerate}

\section{Autograder: the benefit does not transfer}

The graded programs are compiled binaries whose sources are unavailable, so everything
through section~8 rests on the assumption that they contain the pattern and contain it often.
Three submissions settle the question, differing in one feature at a time:

\begin{center}
\small
\begin{tabular}{lrrrr}
\toprule
graded program & October & + prediction & + fusion v1 & + fusion v2 \\
\midrule
Perf        & 0.587 & 0.499 & 0.503 & 0.514 \\
Fibonacci   & 0.586 & 0.499 & 0.500 & 0.499 \\
Bubble sort & 0.248 & 0.236 & 0.236 & 0.236 \\
Check vowel & 0.183 & 0.174 & 0.175 & 0.174 \\
Factorial   & 0.147 & 0.127 & 0.128 & 0.127 \\
GCD         & 0.084 & 0.069 & 0.069 & 0.068 \\
Sum array   & 0.057 & 0.054 & 0.053 & 0.052 \\
Swap shift  & 0.017 & 0.016 & 0.017 & 0.014 \\
Swap        & 0.016 & 0.016 & 0.015 & 0.015 \\
Simple if   & 0.013 & 0.012 & 0.012 & 0.012 \\
Simple add  & 0.012 & 0.011 & 0.011 & 0.011 \\
\midrule
total (s)   & 1.950 & 1.713 & 1.719 & 1.722 \\
vs prediction & -- & $-12.2$\% & $+0.4$\% & $+0.5$\% \\
\bottomrule
\end{tabular}
\end{center}

All forty-four runs report CORRECT. Prediction is the branch-history table and BTB;
v1 adds this chapter's original rule; v2 adds the mirrored operand form.

\textbf{Prediction: $-12.2$\% of simulated time. Fusion v1: $+0.4$\%. Fusion v2: $+0.5$\%,
with Perf 3\% slower.} Both fusion variants are indistinguishable from noise, and nominally
worse. All correct: the fused operation computes the right answer on every graded program
in both forms. It simply almost never fires.

\subsection{Why the transfer failed, and why v2 did not rescue it}

The measurement is sensitive enough to see the effect being tested. Prediction removes about
three cycles per loop iteration and shows up as 12--18\% on the loop-heavy programs
(Perf $-15.0$\%, GCD $-17.9$\%, Fibonacci $-14.8$\%). Fusion removes one instruction per
fused pair and should have been visible at half that magnitude. Neither variant is visible at
all, so the pairs are not being fused.

The v1 hypothesis was that the rule's operand requirement was too narrow:

\begin{verbatim}
    addi x5, x5, -1
    blt  x5, x6, loop     rs1 == rd   ->  v1 fuses
    bge  x6, x5, loop     rs2 == rd   ->  v2 fuses   (v1 rejects)
\end{verbatim}

v2 accepts both forms, so between them the two variants cover every adjacent
\texttt{addi}+branch in which the branch tests the addi's result. \textbf{The union still shows
nothing}, which refutes the operand-order explanation: the graded binaries either contain
almost none of these pairs in either form, or the saving is small enough to be buried.

The local result stands on its own terms --- 1339 to 1035 cycles, $-22.7$\% --- and it is
measured on code that contains the pattern roughly once per loop iteration. What the
submissions show is that the graded programs are not that code. Two candidate reasons remain,
and the data does not separate them:

\begin{itemize}
\item \textbf{The pattern is genuinely rare in compiled loops.} Loop counters are often kept
in memory or updated with other instruction shapes, and a compiler may test a pointer or an
index derived elsewhere rather than the register the \texttt{addi} just wrote.
\item \textbf{The autograder's metric blunts the effect.} It reports simulation wall time, and
the fusion build adds a second combinational instruction-memory read port. That costs
simulator time per cycle regardless of architectural cycles, so a real cycle saving can be
partly cancelled. Against this, the two heaviest loop programs are unchanged to the
millisecond while Perf is 3\% \emph{worse}, which argues the cycle effect itself is small --
but it cannot be separated from the simulation cost with the numbers the grader reports.
\end{itemize}

\subsection{What follows}

\begin{enumerate}
\item \textbf{Neither variant should be shipped.} Both are correct, both add a compare and a
32-bit adder to the fetch path, and both return approximately zero on the only programs that
are graded. Per the project's own decision rule --- adopt only after the log shows a benefit
--- the answer is \texttt{FUSE\_ENABLE = 0}, and the branch-prediction build is the one to
submit.
\item \textbf{The variant space is now exhausted cheaply.} v2 was the only extension the
register file permits: a general "any adjacent \texttt{addi}+branch" rule needs three source
reads and the file has two ports. So this is not "fusion was tried once and abandoned" --- the
reachable design space was enumerated, two points were built, measured and rejected.
\item \textbf{The $-22.7$\% in sections 3--8 stands, with its scope stated.} It is the saving
on code that contains a fused-form pair: one cycle per pair, exactly as predicted. It is not,
and must not be quoted as, a saving on the graded programs.
\item \textbf{Effort belongs on prediction.} The same three submissions show prediction
returning $-12.2$\% on the graded suite, roughly thirty times the fusion variants on the
programs that are actually scored.
\end{enumerate}

\section{Epilogue: the optimisation that replaced it}

Everything above is a negative result with a diagnosis attached. The diagnosis said the
graded programs are dominated by mispredicts, not by repeated instructions; this section is
what that implies, and it is the part that paid.

\subsection{The mechanism}

A mispredict resolved in execute squashes two instructions. Resolved one stage earlier, in
decode, it squashes one. Conditional branches therefore compare in D:

\begin{verbatim}
    d_npc        = d_brtaken ? d_pc + d_imm : d_pc + 4
    mispredict_d = (d_opcode == B_TYPE) && !hazard && (fd_pred_npc != d_npc)
\end{verbatim}
Three things make that non-trivial, and all three cost a bug during bring-up:

\begin{itemize}
\item \textbf{Operands are still in flight.} D forwards from X and M. A load in M must supply
the memory output, not the ALU result, because \texttt{xm\_wdata} has no memory-data case ---
which is exactly why the existing M$\rightarrow$X forwarding masks loads out and lets them
reach W first. D has no such option.
\item \textbf{The X-stage value is not always the ALU result.} \texttt{lui}/\texttt{auipc}
write the immediate and \texttt{jal} writes $\mathrm{pc}+4$, selected only at writeback. A
branch comparing against a \texttt{lui} result read zero.
\item \textbf{The X-stage check must be removed.} Leaving \texttt{B\_TYPE} in the execute-stage
prediction check makes the X redirect win over the D redirect for the same branch: correct,
67/67 passing, and exactly zero cycles saved.
\end{itemize}

\subsection{Local result}

\begin{center}
\small
\begin{tabular}{lrrrr}
\toprule
kernel & baseline cycles & early resolution & saved & squashes before/after \\
\midrule
\texttt{sweep4}  &  29 &  26 & $-3$ & 3/0 \\
\texttt{sweep10} &  47 &  44 & $-3$ & 3/0 \\
\texttt{sweep25} &  92 &  89 & $-3$ & 3/0 \\
\texttt{sweep50} & 167 & 164 & $-3$ & 3/0 \\
\texttt{sweep100}& 317 & 314 & $-3$ & 3/0 \\
\texttt{nested}  & 112 & 102 & $-10$ & 10/0 \\
\texttt{if-in-loop}& 56 &  50 & $-6$ & 6/0 \\
\texttt{smul}    & 218 & 213 & $-5$ & 5/0 \\
\texttt{fib}     &  69 &  66 & $-3$ & 3/0 \\
\midrule
\textbf{total, 14 kernels} & \textbf{1339} & \textbf{1288} & $\mathbf{-51\ (-3.8\%)}$ & \\
\bottomrule
\end{tabular}
\end{center}

The saving is exactly one cycle per mispredict on every kernel: \texttt{nested} saves ten for
its ten, the loop kernels save three each --- entry, exit, and the cold first iteration. All
67 tests pass.

\subsection{Graded result}

\begin{center}
\small
\begin{tabular}{lrrrrc}
\toprule
graded program & October & + prediction & + early resolution & + load-use \& jalr & vs previous \\
\midrule
Perf        & 0.587 & 0.499 & 0.487 & \textbf{0.453} & $-7.0$\% \\
Fibonacci   & 0.586 & 0.499 & 0.489 & \textbf{0.474} & $-3.1$\% \\
Bubble sort & 0.248 & 0.236 & 0.231 & \textbf{0.196} & $-15.2$\% \\
Factorial   & 0.147 & 0.127 & 0.123 & \textbf{0.113} & $-8.1$\% \\
GCD         & 0.084 & 0.069 & 0.068 & \textbf{0.065} & $-4.4$\% \\
Sum array   & 0.057 & 0.054 & 0.052 & \textbf{0.045} & $-13.5$\% \\
Check vowel & 0.183 & 0.174 & 0.173 & 0.175 & $+1.2$\% \\
Swap        & 0.016 & 0.016 & 0.015 & 0.014 & $-6.7$\% \\
Swap shift  & 0.017 & 0.016 & 0.016 & 0.015 & $-6.3$\% \\
Simple if   & 0.013 & 0.012 & 0.012 & 0.012 & $0.0$\% \\
Simple add  & 0.012 & 0.011 & 0.011 & 0.011 & $0.0$\% \\
\midrule
total (s)   & 1.950 & 1.713 & 1.677 & \textbf{1.573} & $-6.2$\% \\
vs October  & --    & $-12.2$\% & $-14.0$\% & $\mathbf{-19.3\%}$ & \\
\bottomrule
\end{tabular}
\end{center}

All eleven CORRECT. Between the two early-resolution builds, eight programs are faster, two
unchanged and one $+1.2$\% --- $2\,\mathrm{ms}$, which is the grader's quantum, not a
regression.

\subsection{The local suite was wrong by a factor of six}

The load-use removal was worth $1.0$\% on the project's fourteen kernels and **$6.2$\%** on the
graded programs. The mechanism is identical --- one cycle per load-use pair --- so the
difference is the workload: the hand-written kernels are pure arithmetic loops with a handful
of loads, while bubble sort and sum array spend their inner loops loading array elements, which
is why they move by $15.2$\% and $13.5$\%.

The contrast with fusion is the point of this whole document:

\begin{center}
\begin{tabular}{lll}
\toprule
feature & local suite said & graded programs said \\
\midrule
macro-op fusion & $-22.7$\% & $+0.4$\% (overstated by $\approx 50\times$) \\
early resolution & $-4.8$\% & $-6.2$\% (understated by $\approx 1.3\times$) \\
\bottomrule
\end{tabular}
\end{center}

A hand-written kernel suite can be wrong in both directions, and only the graded programs say
which. That is the argument for the autograder comparison being part of the experiment loop
rather than an afterthought.

\subsection{The feature ledger}

\begin{center}
\begin{tabular}{llrr}
\toprule
change & mechanism & graded effect & local \\
\midrule
pre-decoder + BHT/BTB & remove mispredicts & $-12.2$\% & $-20$\% \\
macro-op fusion v1    & remove instructions & $+0.4$\% & $-21.7$\% \\
macro-op fusion v2    & remove instructions & $+0.5$\% & $-22.7$\% \\
early resolution      & make mispredicts cheaper & $-2.1$\% & $-3.8$\% \\
load-use + jalr in D  & remove stall cycles & $\mathbf{-6.2}$\% & $-1.0$\% \\
\midrule
cumulative            & & $\mathbf{-19.3}$\% (1.950 $\rightarrow$ 1.573\,s) & \\
\bottomrule
\end{tabular}
\end{center}

Two features measured and rejected, two measured and kept, each with the same three-build
discipline. The rejection of fusion is what pointed at mispredicts; the mispredict work is what
made the pipeline fast enough that load-use stalls became visible. Each step was found by
measuring the previous one on the programs that are actually scored.

\section{What this measurement does not show}

\begin{itemize}
\item \textbf{Area and timing.} Neither is measured. The fetch-path adder is the part most
likely to move $f_{max}$, and section~8.3 states the budget it has to stay inside
($f_{on}/f_{off} > 0.782$) rather than a result. Synthesis on sky130 is still outstanding.
\item \textbf{Any benefit on the graded programs.} Fusion was measured there and returns
$+0.4$\%, which is noise (section~9). The 21.7\% applies to code containing the v1 operand
pattern, not to the graded binaries.
\item \textbf{Runtime configurability.} \texttt{FUSE\_ENABLE} is a compile-time constant, not
a CSR or a pin. Switching it requires a rebuild.
\item \textbf{Other fusion variants.} The measured ranking put \texttt{addi}+\texttt{branch}
at 21.4\% of cycles and every memory-side candidate at 0.9\% or below. Nothing here says
anything about those.
\end{itemize}

\section{Reproducing}

\begin{verbatim}
  make sim-all -C verif/          # with FUSE_ENABLE = 1'b1 in code/rv_core.sv
  make sim-all -C verif/          # with FUSE_ENABLE = 1'b0, for the other column
  make tb-fuse_detect_tb -C verif/  # the recogniser's own unit TB, 12/12
\end{verbatim}

The unit testbench checks the recogniser against twelve rows -- seven rejections from the
contract's truth table and five acceptances, including a second base address and a pair whose
source and destination registers differ, so that neither \texttt{pc\_i} nor the register
fields can go unread without a failure.


\end{document}
